% The M05 code sheet: a paper handout, separate from the pen-and-paper exercise.
% A student types each framed box into the matching empty cell of a Colab
% notebook and looks at what comes out. The sheet is about how to use the tools.
%
% Compile with: xelatex -interaction=nonstopmode code-sheet.tex   (run twice)
%
% The boxes (boxes/box*.tex) are GENERATED, together with the notebooks, from
% one list in tools/build_m05_code_sheet.py, so that the code on the paper is
% the code that was run. Edit the list, not the files:
%
%     python tools/build_m05_code_sheet.py
\documentclass[a4paper, 12pt]{extarticle}
\usepackage[top=0.7in, bottom=0.65in, left=0.75in, right=0.75in]{geometry}
\usepackage{amsmath}
\usepackage{graphicx}
\usepackage[hidelinks]{hyperref}   % hidelinks, or the print carries colour boxes
\usepackage{enumitem}
\usepackage{listings}
\usepackage{tcolorbox}
\tcbuselibrary{listings}
\usepackage{fontspec}
\usepackage{xcolor}
\usepackage{tikz}

% Body face: Charter, asked for the way exercise.tex asks for it.
\IfFontExistsTF{XCharter-Roman.otf}{
  \setmainfont{XCharter}[
    Extension      = .otf,
    UprightFont    = *-Roman,
    BoldFont       = *-Bold,
    ItalicFont     = *-Italic,
    BoldItalicFont = *-BoldItalic]
}{
  \IfFontExistsTF{Charter}{\setmainfont{Charter}}{
    \IfFontExistsTF{Iowan Old Style}{\setmainfont{Iowan Old Style}}{
      \IfFontExistsTF{Georgia}{\setmainfont{Georgia}}{}}}}

% Code face: one that tells 0 from O and l from 1, because it is copied by hand.
% A chain, not a name: naming a font that is not installed is a hard error.
\IfFontExistsTF{JetBrains Mono}{\setmonofont{JetBrains Mono}[Scale=0.92]}{
  \IfFontExistsTF{DejaVu Sans Mono}{\setmonofont{DejaVu Sans Mono}[Scale=0.9]}{
    \IfFontExistsTF{Menlo}{\setmonofont{Menlo}[Scale=0.92]}{}}}

\setlength{\parindent}{0pt}
\setlength{\parskip}{0.4em}
\setlength{\emergencystretch}{1.2em}
\raggedbottom
% No single line of a note is left alone at the foot or the head of a page.
\clubpenalty=10000
\widowpenalty=10000

% Break the page unless #1 of it is left (needspace's glue trick, written out
% because needspace.sty is not in a BasicTeX install; see exercise.tex).
\newcommand{\needroom}[1]{%
  \par
  \begingroup
    \setlength{\dimen0}{#1}%
    \vskip 0pt plus \dimen0
    \penalty -100%
    \vskip 0pt plus -\dimen0
  \endgroup}

% A part heading. It may not be the last thing on a page, but it is not made to
% pull a whole part after it: a part is a heading, a box and its notes.
\newcommand{\parthead}[2]{%
  \needroom{5\baselineskip}%
  \vspace{0.5em}
  {\large\bfseries Part #1 --- #2}\par
  \vspace{-0.4em}\rule{\textwidth}{0.8pt}\par
  \vspace{0.1em}}

% The framed code. The student types what is inside, and only that. Line numbers
% are in the margin so that the notes under the box can point at a line.
\lstset{
  basicstyle=\ttfamily\fontsize{10.5}{14}\selectfont,
  columns=fixed, keepspaces=true, showstringspaces=false, upquote=true,
  numbers=left, numberstyle=\ttfamily\tiny\color{black!55}, numbersep=7pt,
  xleftmargin=1.4em, tabsize=4,
}
\newtcblisting{codebox}[2]{
  listing only, colback=white, colframe=black, boxrule=1.1pt, arc=2pt,
  top=2pt, bottom=2pt, left=2pt, right=2pt, boxsep=2pt,
  title={\bfseries #1\hspace{0.9em}\mdseries #2},
  coltitle=black, colbacktitle=black!9, fonttitle=\normalsize,
  before skip=0.5em, after skip=0.1em, nobeforeafter,
}
\newenvironment{linenotes}{%
  \small
  \begin{description}[font=\bfseries\ttfamily\footnotesize, leftmargin=2.9em,
    labelwidth=2.4em, labelsep=0.3em, itemsep=0.12em, parsep=0pt, topsep=0.25em]}%
  {\end{description}}

% The play button of a Colab cell, drawn, because no text font here has the glyph.
\newcommand{\playbutton}{\tikz[baseline=-0.55ex]{\draw[fill=black!80] (0,0) circle (1.05ex);\fill[white] (-0.3ex,-0.55ex) -- (0.6ex,0) -- (-0.3ex,0.55ex) -- cycle;}}

% The lab. Both the QR code and the printed address are the course's own short
% link, never the notebook's: paper outlives any one upload (see exercise.tex).
\newcommand{\labaddress}{go.skojaku.com/m05lab}
\newcommand{\labtarget}{https://go.skojaku.com/m05lab}

\begin{document}

{\LARGE\bfseries Lab \textperiodcentered{} Find the blocks}\hfill{\large the code sheet}

\vspace{0.3em}
\begin{tcolorbox}[colback=white, boxsep=1pt, top=4pt, bottom=4pt]
\begin{minipage}[c]{0.2\textwidth}
\begin{center}
\includegraphics[width=2.6cm]{m05lab-qr.png}
\end{center}
\end{minipage}
\hfill
\begin{minipage}[c]{0.77\textwidth}
\textbf{Before you start.} Point a camera at the square, or type
\href{\labtarget}{\texttt{\labaddress}}. It opens a notebook in Google Colab.
Press the \playbutton{} button on each of its three \emph{setup} cells, in order. The first
one restarts Colab by itself, which is meant to happen. The second takes a few
minutes. Wait for the \texttt{\char"2713} under the third. You type nothing in the
setup cells.
\end{minipage}
\end{tcolorbox}

\textbf{How to use this sheet.} Each framed \textbf{Box} below goes into the
matching empty cell of the notebook. \textbf{Type it by hand} rather than pasting it:
writing a line is how you find out what it does. Press \textbf{Shift+Enter} to run
a cell. The notes under a box say what each numbered line does.

\textbf{Typing tips.} Spelling, capital letters, brackets and quotation marks all
matter. Colab may add a closing bracket or quote for you; when you type the closing
one yourself, it steps over the extra one. Indent with the Tab key.

\vspace{0.2em}
\begin{tcolorbox}[colback=white, boxsep=1pt, top=3pt, bottom=3pt, title={\bfseries If something goes wrong}, coltitle=black, colbacktitle=black!9, fonttitle=\normalsize]
\footnotesize
\begin{itemize}[itemsep=0.05em, topsep=0.05em, leftmargin=1.2em]
\item A red message ending in \texttt{NameError} means a name does not exist yet:
you skipped an earlier box, or Colab was reset. Run the earlier boxes again, in order.
\item \texttt{SyntaxError} means a typing slip. Look for a missing bracket or quote
near the \texttt{\^{}} that the message points at.
\item The \texttt{\char"2713} is gone, or Colab says it has disconnected: Colab forgets
everything after a break. Run the three setup cells again. Anything else: raise your hand.
\end{itemize}
\end{tcolorbox}

\parthead{1}{Meet the data}

Zachary's karate club: 34 members of a university karate club, studied in the 1970s,
linked when they were friends outside class. The club split in two after a quarrel
between the instructor, \textbf{Mr.\ Hi} (node 0), and the administrator, \textbf{the
Officer} (node 33). We know who went with whom, so it is the standard test case.

% GENERATED by tools/build_m05_code_sheet.py -- do not edit.
\begin{codebox}{Box 1}{Build the karate club network}
import igraph as ig
import graph_tool.all as gt
karate = ig.Graph.Famous("Zachary")
g = gt.Graph(directed=False)
g.add_edge_list(karate.get_edgelist())
print(g.num_vertices(), g.num_edges())
\end{codebox}
\begin{linenotes}
\item[1] \texttt{import} loads a library: a bundle of ready-made tools. \texttt{igraph} is a library for networks. We use it for one thing only: Leiden. \texttt{as ig} is its short name.
\item[2] \texttt{graph\_tool} is the main library of this sheet. It fits stochastic block models and draws networks. \texttt{as gt} is its short name.
\item[3] \texttt{ig.Graph.Famous("Zachary")} builds a network that igraph already knows: Zachary's karate club. The \texttt{=} stores it under the name \texttt{karate}.
\item[4--5] graph-tool cannot read igraph's networks. It has its own kind, so we make an empty one (\texttt{directed=False} says that a friendship goes both ways) and fill it. \texttt{karate.get\_edgelist()} asks igraph for all the friendships as pairs of node numbers, and \texttt{add\_edge\_\allowbreak list} hands the pairs to graph-tool, which creates the nodes and the edges. This is the only bridge between the two libraries.
\item[6] \texttt{num\_vertices()} counts the nodes (the members) and \texttt{num\_edges()} counts the edges (the friendships). You should see \texttt{34 78}.
\end{linenotes}


\parthead{2}{Find communities with Leiden}

% GENERATED by tools/build_m05_code_sheet.py -- do not edit.
\begin{codebox}{Box 2}{Find communities with Leiden}
result = karate.community_leiden("modularity", n_iterations=-1)
leiden = result.membership
print(leiden)
\end{codebox}
\begin{linenotes}
\item[1] \texttt{community\_leiden} runs the Leiden algorithm. It looks for \emph{communities}: groups with many friendships inside and few between groups. \texttt{"modularity"} is the score it tries to raise: how much denser the groups are than in a random network. \texttt{n\_iterations=-1} means keep improving until nothing changes. What it returns is stored as \texttt{result}.
\item[2] \texttt{result.membership} is the answer as a plain list: entry $i$ is the group number of node $i$. We store it as \texttt{leiden}.
\item[3] Print the list. You should see 34 numbers. Two nodes with the same number are in the same group. The numbers are only labels, and Leiden uses chance, so your list may differ from your neighbour's.
\end{linenotes}


\parthead{3}{Draw them with graph-tool}

From here on, graph-tool does the drawing.

% GENERATED by tools/build_m05_code_sheet.py -- do not edit.
\begin{codebox}{Box 3}{The stochastic block model, with graph-tool}
import graph_tool.all as gt
import numpy as np
gt.seed_rng(1)
g = gt.Graph(directed=False)
g.add_edge_list(karate.get_edgelist())
state = gt.minimize_blockmodel_dl(g, state_args=dict(deg_corr=False))
sbm = state.get_blocks().a
print(sbm)
groups = np.unique(sbm, return_inverse=True)[1]
colors = g.new_vertex_property("vector<double>")
for v in g.vertices():
    colors[v] = (*palette[groups[int(v)]], 1)
pos = gt.graph_draw(
    g,
    vertex_fill_color=colors,
    vertex_text=g.vertex_index,
)
\end{codebox}
\begin{linenotes}
\item[1] graph-tool is a second library. It fits the stochastic block model, which you know from the lecture. \texttt{as gt} is its short name.
\item[2] \texttt{numpy} is a library for lists of numbers. \texttt{np} is its short name.
\item[3] The fit makes random choices. \texttt{seed\_rng(1)} fixes them, so that everyone in the room gets the same answer. Change the \texttt{1} and you may get a different grouping.
\item[4] graph-tool cannot read igraph's networks. It has its own kind, so we start with an empty one. \texttt{directed=False} says that a friendship goes both ways.
\item[5] \texttt{karate.get\_edgelist()} asks igraph for all the friendships as pairs of node numbers. \texttt{add\_edge\_\allowbreak list} hands the pairs to graph-tool, which creates the nodes and the edges. This line is the only bridge between the two libraries.
\item[6] \texttt{minimize\_blockmodel\_dl} fits the stochastic block model. \texttt{dl} stands for \emph{description length}: of all the ways to group the nodes, it picks the one that describes the network in the fewest bits. It decides the number of groups too. \texttt{state\_args=dict(deg\_corr=False)} asks for the plain model of the lecture, where the group alone decides how likely a friendship is. Graph-tool's default also lets every member have a number of friends of their own; on a network this small, that version finds one single group. The fitted model is stored as \texttt{state}.
\item[7] \texttt{state.get\_blocks()} gives the group of each node, in graph-tool's own array type. \texttt{.a} turns it into a plain array of numbers, like the list in Box 2.
\item[8] Print it: again one number for each of the 34 nodes.
\item[9] graph-tool numbers its groups with gaps (say 15 and 32). \texttt{np.unique} with \texttt{return\_\allowbreak inverse=True} renumbers them 0, 1, 2, \ldots{} so that they fit the palette; the \texttt{[1]} picks that renumbered list.
\item[10] graph-tool keeps what it knows about nodes in a \emph{property map}: a table with one value per node. This one will hold a list (red, green, blue) for each node.
\item[11--12] A loop over the nodes \texttt{v} of the network. \texttt{groups[int(v)]} is the group of node \texttt{v} (\texttt{int} turns graph-tool's node object into its number), \texttt{palette[\ldots]} is the colour of that group, and \texttt{(*palette[\ldots], 1)} unpacks its red, green and blue and adds a 1 for fully opaque. The result goes into the table.
\item[13--17] \texttt{gt.graph\_draw} is graph-tool's drawing function. \texttt{vertex\_fill\_color=colors} fills each node with its colour from the table, and \texttt{vertex\_text=g.vertex\_index} writes its number on it. The function hands back the position it gave each node; we keep that as \texttt{pos}, which also stops Colab from printing it under the picture.
\end{linenotes}


\parthead{4}{The stochastic block model}

% GENERATED by tools/build_m05_code_sheet.py -- do not edit.
\begin{codebox}{Box 4}{The stochastic block model}
deg_corr = False
gt.seed_rng(1)
state = gt.minimize_blockmodel_dl(
    g, state_args=dict(deg_corr=deg_corr)
)
sbm = state.get_blocks().a
print(sbm)
pos = gt.graph_draw(
    g, pos=pos, vertex_text=g.vertex_index,
    vertex_fill_color=paint(sbm),
)
\end{codebox}
\begin{linenotes}
\item[1] This stores your choice of model under the name \texttt{deg\_corr}. \texttt{False} is the plain model of the lecture: the group alone decides how likely a friendship is, so the members of a group are alike in how many friends they have. \texttt{True} is the \emph{degree-corrected} model: every member may also have more or fewer friends than the others in the group. Change this one word later, to try the other model.
\item[2] The fit makes random choices. \texttt{seed\_rng(1)} fixes them, so that everyone in the room gets the same answer. Change the \texttt{1} and you may get a different grouping.
\item[3--5] \texttt{minimize\_blockmodel\_dl} fits the stochastic block model from the lecture. \texttt{dl} stands for \emph{description length}: of all the ways to group the nodes, it picks the one that describes the network in the fewest bits. It decides the number of groups too. \texttt{state\_args=dict(deg\_corr=deg\_corr)} hands your choice from line 1 to the model. The fitted model is stored as \texttt{state}.
\item[6] \texttt{state.get\_blocks()} gives the group of each node, in graph-tool's own array type. \texttt{.a} turns it into a plain array of numbers, like the list in Box 2.
\item[7] Print it: one number for each of the 34 nodes.
\item[8--11] The drawing call of Box 3, with the colours of the block model's groups. \texttt{pos=pos} reuses the positions from Box 3, so that each node stays in the same place and the pictures can be compared.
\end{linenotes}


\textbf{Try it.} Change \texttt{False} to \texttt{True} in line 1 and run the cell
again. How many groups does the model find now, and what happens to the picture?
Then change it back to \texttt{False}.

\parthead{5}{What really happened}

% GENERATED by tools/build_m05_code_sheet.py -- do not edit.
\begin{codebox}{Box 5}{Draw the groups on a circle}
import numpy as np
order = np.argsort(leiden)
layout = karate.layout_circle(order=order)
fig, ax = plt.subplots()
ig.plot(result, target=ax, layout=layout, vertex_label=karate.vs.indices)
plt.show()
\end{codebox}
\begin{linenotes}
\item[1] \texttt{numpy} is a library for lists of numbers. \texttt{np} is its short name.
\item[2] \texttt{np.argsort(leiden)} lists the node numbers sorted by group, so that nodes of one group sit together.
\item[3] \texttt{layout\_circle(order=order)} puts the nodes on a circle in that order, so that each group takes one arc. The layout is a position for every node.
\item[4] A fresh empty panel.
\item[5] The drawing call again, with \texttt{layout=layout} to use our positions instead of igraph's own.
\item[6] Show the picture.
\end{linenotes}


\textbf{What to look at.} Put your three pictures side by side (the block model with
\texttt{deg\_corr = False}): the nodes are in the same places. Find node 0 and node 33. Which groupings keep them apart? In the
stochastic block model's picture, who ends up together, and how many friends do they have?

\parthead{6}{A hierarchy of groups}

% GENERATED by tools/build_m05_code_sheet.py -- do not edit.
\begin{codebox}{Box 6}{The US airports: Leiden, and a map}
import seaborn as sns
import pandas as pd
members = g_air.community_leiden("modularity", n_iterations=-1).membership
top = pd.Series(members).value_counts().index[:4]
sns.scatterplot(x=lon, y=lat, color="#bdbdbd", s=14)
sns.scatterplot(
    x=lon, y=lat, hue=members, hue_order=top, palette=palette[:4], s=14
)
plt.show()
\end{codebox}
\begin{linenotes}
\item[1--2] \texttt{seaborn} draws statistical pictures, and \texttt{pandas} works with tables. \texttt{sns} and \texttt{pd} are their usual short names.
\item[3] Box 2, line 1, again, on the US airport network \texttt{g\_air} (540 airports, two linked when a flight connects them), which is ready-made. We keep the group of each airport as \texttt{members}.
\item[4] \texttt{pd.Series(members).value\_counts()} counts how many airports each group has, biggest group first. \texttt{.index[:4]} keeps the numbers of the four biggest groups. We call them \texttt{top}.
\item[5] \texttt{lon} and \texttt{lat} are the longitude and the latitude of each airport, also ready-made. \texttt{sns.\allowbreak scatterplot} puts one dot per airport at that place: a rough map of the US. This first call paints all of them light grey; \texttt{s=14} is the dot size.
\item[6--8] The second call paints over the grey, but only the four groups in \texttt{top}: \texttt{hue=members} says that the colour follows the group, \texttt{hue\_order=top} limits it to those four, and \texttt{palette[:4]} is the first four colours of the palette. Every other airport stays light grey: it is ``the rest''.
\item[9] Show the map.
\end{linenotes}


\textbf{What to look at.} How many levels does the hierarchy have, and how many groups
at each? Now change \texttt{False} to \texttt{True} and run the cell again. What changes?

\parthead{7}{If you have time: the US airports}

The same recipe on a bigger network: 540 US airports, two linked when a flight
connects them.

% GENERATED by tools/build_m05_code_sheet.py -- do not edit.
\begin{codebox}{Box 7}{The US airports: the hierarchy}
deg_corr = False
g2 = gt.Graph(directed=False)
g2.add_edge_list(g_air.get_edgelist())
nested_air = gt.minimize_nested_blockmodel_dl(
    g2, state_args=dict(deg_corr=deg_corr)
)
nested_air.print_summary()
pos_air = nested_air.draw()
\end{codebox}
\begin{linenotes}
\item[1] Your choice of model, as in Box 4.
\item[2--3] Box 1, lines 4 and 5, again, with a new graph-tool network called \texttt{g2}. \texttt{g\_air} is the US airport network (540 airports, two linked when a flight connects them), already built for you.
\item[4--6] Box 6, lines 2 to 4, again, on the airport network.
\item[7] One line for each level of the hierarchy, as in Box 6.
\item[8] The hierarchy on a circle. Without \texttt{vertex\_fill\_color}, graph-tool picks the colours itself; with this many groups, colour no longer says much, and the place on the circle does.
\end{linenotes}


\textbf{Try it.} Run the cell with \texttt{False}, then with \texttt{True}. Compare the
printed levels and the two pictures. The number of flights differs a lot from one
airport to another. Which setting describes this network better, do you think, and why?

\textbf{Then try your own.} Change the network in Box 1, and run Boxes 2 to 4 and
Box 6 again. A small network to try instead of the karate club:

\hspace*{1.4em}\texttt{ig.Graph.Famous("Krackhardt\_Kite")}

\end{document}
